\chapter{Contracepção}\label{ch:g8:contraception}

Os últimos capítulos mapearam um mecanismo de poder notável: duas
\omterm{def:g5:first-look-at-cells:cell}{células} se encontram e, quarenta semanas depois, existe uma pessoa. Os
seres humanos são a única \omterm{def:g5:biodiversity:species}{espécie} que entende o próprio mecanismo --- e
entender traz uma escolha que o sapo nunca enfrenta: se e quando pôr
esse mecanismo para funcionar. A \emph{\omterm{prop:g8:contraception:principle}{contracepção}} é a ciência
aplicada dessa escolha, e ela funciona, em todos os métodos dela,
interrompendo um passo que você agora sabe nomear.

\section{O princípio: nomear o elo, bloquear o elo}

\begin{proposition}[A contracepção se encaixa no mecanismo]\label{prop:g8:contraception:principle}
A \omterm{def:g5:human-reproduction:pregnancy}{gravidez} exige uma corrente ininterrupta
(\cref{prop:g8:sexual-reproduction:core},
\cref{prop:g8:from-egg-to-birth:firstweek}): \omterm{def:g8:sexual-reproduction:gametes}{gametas} produzidos ---
\omterm{def:g8:reproductive-systems:male}{espermatozoides} depositados --- a nadada até a \omterm{def:g8:reproductive-systems:female}{tuba uterina} --- uma
\omterm{def:g4:animal-reproduction:sexual}{célula reprodutora feminina} liberada e encontrada --- \omterm{def:g5:flower-to-fruit:steps}{fecundação} ---
implantação. A \emph{contracepção}\index{contracepção} é qualquer
interrupção deliberada dessa corrente. Todo método da prateleira da
farmácia é um ataque a um elo com nome, e classificá-los pelo elo
deles é a teoria inteira:
\begin{itemize}
  \item os métodos de \emph{barreira} impedem os \omterm{def:g8:sexual-reproduction:gametes}{gametas} de se
        encontrarem;
  \item os métodos \emph{hormonais} impedem que a \omterm{def:g4:animal-reproduction:sexual}{célula reprodutora
        feminina} seja liberada;
  \item outros métodos fazem o encontro ou a implantação falharem.
\end{itemize}
\end{proposition}
\admitted

\section{Os principais métodos, pelo elo de cada um}

\begin{definition}[Métodos de barreira: a camisinha]\label{def:g8:contraception:condom}
A \emph{camisinha}\index{camisinha} --- uma capa fina usada no \omterm{def:g8:reproductive-systems:male}{pênis}
(uma versão feminina forra a vagina) --- bloqueia a corrente no passo
do depósito: o sêmen nunca entra no corpo do parceiro. É de fácil
acesso e só é necessária quando é necessária --- e é única na
prateleira por um segundo serviço: é o único método que também bloqueia
a passagem das \emph{infecções} transmitidas pelo sexo, o mundo do
\cref{ch:g9:microbes-and-infection}. É esse serviço duplo que mantém a
camisinha em cena mesmo quando outro método está em uso.
\end{definition}

\begin{definition}[Métodos hormonais: a pílula]\label{def:g8:contraception:pill}
A \emph{pílula anticoncepcional}\index{pílula anticoncepcional} usa a
própria maquinaria da \cref{def:g8:puberty:hormones}: tomada todo dia,
os \omterm{def:g8:puberty:hormones}{hormônios} dela mantêm o ciclo dos \omterm{def:g8:reproductive-systems:female}{ovários} no estado suspenso (o
acontecimento 4 da \cref{prop:g8:reproductive-systems:cycle} sem
\omterm{def:g5:human-reproduction:pregnancy}{gravidez} nenhuma) --- \emph{sem \omterm{prop:g8:reproductive-systems:cycle}{ovulação}, sem \omterm{def:g4:animal-reproduction:sexual}{célula reprodutora
feminina}, sem nada para fecundar}. É receitada por um médico, que
ajusta o método à pessoa; é muito confiável quando tomada com
regularidade --- e não oferece nada contra as infecções, daí a parceria
permanente com a \omterm{def:g8:contraception:condom}{camisinha}.
\end{definition}

\begin{example}[A prateleira, classificada por elo]\label{ex:g8:contraception:shelf}
Os demais métodos comuns, arquivados pelo elo que atacam: os implantes
e adesivos hormonais --- o princípio da pílula com memória melhor
(meses ou anos, sem tomar nada todo dia); os pequenos dispositivos que
um médico coloca no \omterm{prop:g5:human-reproduction:plan}{útero} --- fazem o encontro e a implantação falharem
por anos a fio; e a \emph{pílula do dia seguinte}, uma dose hormonal
depois do fato, dentro de poucos dias, que funciona principalmente
atrasando uma \omterm{prop:g8:reproductive-systems:cycle}{ovulação} que ainda não aconteceu --- um freio de último
recurso, disponível na farmácia, e nenhum substituto para um método
usado \emph{antes}. O que nenhum deles faz: qualquer coisa contra
infecção.
\end{example}

\begin{remark}[O que não funciona]\label{rem:g8:contraception:notwork}
Os métodos do pátio da escola, medidos contra a corrente: ``contar os
dias seguros'' esbarra na variabilidade da
\cref{rem:g8:reproductive-systems:living} --- a data da \omterm{prop:g8:reproductive-systems:cycle}{ovulação} não é
conhecível de antemão, e os \omterm{def:g8:reproductive-systems:male}{espermatozoides} sobrevivem dias esperando
por ela; ``tirar a tempo'' esbarra no fato de que os \omterm{def:g8:reproductive-systems:male}{espermatozoides}
viajam antes do final. Os dois falham exatamente onde a teoria do elo
com nome prevê: eles não bloqueiam nenhum elo de forma confiável.
Lenda não é método.
\end{remark}

\section{Escolha, responsabilidade e o mapa da ajuda}

\begin{proposition}[A escolha é real e compartilhada]\label{prop:g8:contraception:choice}
Fatos que o mecanismo impõe:
\begin{itemize}
  \item a \omterm{def:g5:human-reproduction:pregnancy}{gravidez} precisa de um ato, e não da duração de um
        relacionamento: a corrente não conta ocasiões;
  \item ela é uma corrente \emph{compartilhada}: \omterm{def:g5:first-look-at-cells:cell}{células} dos dois
        parceiros, responsabilidade dos dois --- a \omterm{def:g8:contraception:condom}{camisinha} e a
        pílula não são assunto de um sexo cada, e sim uma conversa;
  \item e nenhum método, tirando a abstinência, é perfeito, embora os
        bons cheguem perto quando usados exatamente como foram
        projetados --- o ``usado direito'' é onde mora quase toda
        falha do mundo real.
\end{itemize}
\end{proposition}
\admitted

\begin{example}[O mapa da ajuda]\label{ex:g8:contraception:help}
Os endereços do \cref{ex:g8:puberty:questions}, ampliados: enfermeiros
escolares e médicos de família respondem a perguntas sobre
\omterm{prop:g8:contraception:principle}{contracepção} com sigilo --- inclusive de menores de idade; existem
centros de planejamento familiar em toda região justamente para isso,
gratuitos e sem julgamento; os farmacêuticos lidam com o relógio da
pílula do dia seguinte. O sistema foi montado para que ninguém precise
navegar a corrente com base em lenda --- o mapa da ajuda faz parte do
método.
\end{example}

\begin{method}[Avaliar qualquer método de que você ouvir falar]\label{met:g8:contraception:evaluate}
Diante de qualquer método alegado, pergunte nesta ordem:
\begin{enumerate}
  \item \emph{que elo} da corrente ele bloqueia --- você consegue
        nomear esse elo no mapa?
  \item \emph{com que confiabilidade}, do jeito que as pessoas reais
        usam?
  \item ele protege também contra as \emph{infecções}, ou a \omterm{def:g8:contraception:condom}{camisinha}
        precisa acompanhá-lo?
  \item o que ele exige --- memória diária, receita, colocação --- e
        isso combina com quem vai usar?
\end{enumerate}
Sem elo nomeável, não há método (a
\cref{rem:g8:contraception:notwork} falha já na pergunta 1).
\end{method}

\section{Exercícios}

\begin{exercise}[$\star$]\label{exo:g8:contraception:1}
Enuncie a corrente que a \omterm{def:g5:human-reproduction:pregnancy}{gravidez} exige e defina \omterm{prop:g8:contraception:principle}{contracepção} em
relação a ela.
\end{exercise}

\begin{exercise}[$\star$]\label{exo:g8:contraception:2}
Que elo a \omterm{def:g8:contraception:condom}{camisinha} bloqueia? E a pílula?
\end{exercise}

\begin{exercise}[$\star$]\label{exo:g8:contraception:3}
Que segundo serviço só a \omterm{def:g8:contraception:condom}{camisinha} oferece?
\end{exercise}

\begin{exercise}[$\star$]\label{exo:g8:contraception:4}
Como a pílula do dia seguinte funciona principalmente, e o que ela não
é?
\end{exercise}

\begin{exercise}[$\star$]\label{exo:g8:contraception:5}
Diga três endereços do mapa da ajuda, com o que cada um oferece.
\end{exercise}

\begin{exercise}[$\star$]\label{exo:g8:contraception:6}
Por que ``contar os dias seguros'' não é confiável? Dois motivos da
\cref{rem:g8:contraception:notwork}.
\end{exercise}

\begin{exercise}[$\star\star$]\label{exo:g8:contraception:7}
Classifique a prateleira do \cref{ex:g8:contraception:shelf} pelo elo
atacado, acrescentando a \omterm{def:g8:contraception:condom}{camisinha} e a pílula.
\end{exercise}

\begin{exercise}[$\star\star$]\label{exo:g8:contraception:8}
Por que os médicos tanto dizem ``pílula \emph{e} \omterm{def:g8:contraception:condom}{camisinha}'' a casais
jovens? Duas proteções, uma frase para cada.
\end{exercise}

\begin{exercise}[$\star\star$]\label{exo:g8:contraception:9}
Explique, com a \cref{prop:g8:reproductive-systems:cycle}, em que
estado a pílula mantém o ciclo e por que isso evita a \omterm{def:g5:human-reproduction:pregnancy}{gravidez}.
\end{exercise}

\begin{exercise}[$\star\star$]\label{exo:g8:contraception:10}
``\omterm{def:g5:first-look-at-cells:cell}{Células} dos dois parceiros, assunto dos dois parceiros.'' Defenda a
frase contra o costume de deixar a \omterm{prop:g8:contraception:principle}{contracepção} para um dos lados do
casal.
\end{exercise}

\begin{exercise}[$\star\star$]\label{exo:g8:contraception:11}
Rode o \cref{met:g8:contraception:evaluate} em: (a) o implante
hormonal; (b) ``tirar a tempo''.
\end{exercise}

\begin{exercise}[$\star\star\star$]\label{exo:g8:contraception:12}
``A \omterm{prop:g8:contraception:principle}{contracepção} é fisiologia aplicada: todo método da prateleira é um
capítulo deste livro, com pontaria.'' Justifique com três métodos e os
três capítulos deles.
\end{exercise}

\section{Problema: A consulta de planejamento familiar}

\begin{problem}[{Problema de fim de semana --- uma tarde de perguntas
no centro de planejamento}]\label{pb:g8:contraception:1}
A tarde (inventada) de uma conselheira de planejamento familiar:
quatro visitantes, quatro situações, um mapa da corrente na parede.

\medskip\noindent\textbf{Parte I --- O mapa da
parede.}
\begin{enumerate}
  \item Desenhe o mapa da parede: a corrente, dos \omterm{def:g8:sexual-reproduction:gametes}{gametas} à
        implantação, com os três pontos de ataque da
        \cref{prop:g8:contraception:principle} marcados.
  \item A visitante 1 pergunta ``como é que uma pílula por dia
        consegue parar alguma coisa?''. Responda com a maquinaria: que
        língua a pílula fala, e com quem?
  \item A visitante 1 emenda: ``e se eu esquecer algumas?''. Situe o
        perigo no mapa --- o que volta a funcionar quando o sinal
        hormonal falha?
  \item Por que a recomendação da \omterm{def:g8:contraception:condom}{camisinha} feita pela conselheira
        sobrevive a qualquer outra escolha de método? (O serviço
        duplo.)
\end{enumerate}

\medskip\noindent\textbf{Parte II --- As situações.}
\begin{enumerate}[resume]
  \item A visitante 2, com dezessete anos e ansiosa, acreditava que
        um aplicativo de dias seguros tornava desnecessários os outros
        métodos. Corrija a crença com gentileza: que dois fatos
        biológicos derrotam o calendário?
  \item A visitante 3 precisa de algo ``sem memória diária''. Ofereça
        os dois candidatos da prateleira e os elos deles.
  \item O visitante 4 chega menos de um dia depois de uma \omterm{def:g8:contraception:condom}{camisinha}
        estourada. Diga a opção, o relógio dela, o mecanismo principal
        dela --- e o método permanente que a conselheira vai discutir
        para depois.
  \item Os quatro saem com a mesma frase final sobre infecções.
        Escreva essa frase.
\end{enumerate}

\medskip\noindent\textbf{Parte III --- A conversa final.}
\begin{enumerate}[resume]
  \item O estagiário pergunta por que o centro ensina a corrente em
        vez de só entregar métodos. Responda com a primeira pergunta
        do \cref{met:g8:contraception:evaluate}.
  \item Classifique os métodos da tarde pelo que cada um exige de
        quem usa (a pergunta 4 do método).
  \item O estagiário se pergunta por que as consultas de menores são
        sigilosas por princípio. Responda com a lógica do
        \cref{ex:g8:contraception:help}.
  \item Encerre a tarde numa frase: o que o centro de fato distribui,
        para além das caixinhas.
\end{enumerate}
\end{problem}
